\section{The Fair Comparison}
\label{sec:fair}

The decision procedure identifies a regime where block-structured decoding
wins on per-event latency. This section asks the sharper question: at a
\emph{fair} comparison point, does that advantage survive? Fairness has two
axes---same quality and same compute---and the pair that satisfies both is the
compute-matched pair of Table~\ref{tab:pair}.

\begin{table}[ht]
\centering
\caption{The compute-matched pair: same quality, comparable active compute.}
\label{tab:pair}
\small
\begin{tabular}{lcc}
\toprule
 & AR-3B & DG-26B-A4B \\
\midrule
Architecture              & dense, autoregressive    & MoE, block diffusion \\
Total parameters          & 3.0B                      & 26B \\
Active parameters         & 3.0B (all)                & 4.0B (top-8 of 128) \\
Active-parameter ratio    & $1.0\times$               & $1.33\times$ \\
Weights (bf16)            & 6.0 GB                    & 52.0 GB \\
9-field F1 ($L{=}256$)    & \arThreeBFnine            & \dgNineF \\
\bottomrule
\end{tabular}
\end{table}

Quality is identical (gap $=0.0002$) and active inference parameters are
comparable ($1.33\times$ ratio). This is the pair for per-event claims. It is
\emph{not} the quality-matched deployment pair: a smaller AR model reaches the
same quality at $17\times$ less memory (Table~\ref{tab:deploy}).

\paragraph{Where DG wins.}
\textit{(a) Per-event latency at batch-1.}
Within a canvas DG emits 1 token or 256 for the same cost; AR pays per token.
The gap widens with output length (Table~\ref{tab:per-event}).

\begin{table}[ht]
\centering
\caption{Per-event latency, batch-1, short input. Fitted values; measured anchors agree within $0.4\%$ (AR) and $1\%$ (DG).}
\label{tab:per-event}
\small
\begin{tabular}{rrrr}
\toprule
$L$ (tokens) & AR-3B (ms) & DG (ms) & DG speedup \\
\midrule
128  & 812   & 379  & $2.14\times$ \\
256  & 1618  & 379  & $4.27\times$ \\
512  & 3228  & 498  & $6.48\times$ \\
1024 & 6448  & 738  & $8.74\times$ \\
2048 & 12889 & 1216 & $10.6\times$ \\
\bottomrule
\end{tabular}
\end{table}

\textit{(b) Output-length invariance within a canvas.}
DG at $L\in\{32,64,128,256\}$ costs $\approx 379$ ms flat (spread $<\!3$ ms);
AR over the same range grows $\approx 209\to 1618$ ms linearly. The canvas
boundary is structural: the model denoises the whole canvas in $K{=}8$ parallel
steps regardless of how many tokens are actually needed.

\textit{(c) Growth rate with output length.}
Each extra canvas adds $\dgStep$ ms to DG; each extra 256 tokens adds
$\arSlope\times\canvasLength\approx 1610$ ms to AR---a $\approx\!13\times$
steeper marginal cost. The per-event advantage therefore \emph{widens} with
output length, from $2.1\times$ at $L=128$ to $10.6\times$ at $L=2048$ (eight
canvases).

\paragraph{Where AR wins.}
\textit{(d) Batched throughput.}
AR's KV cache lets batching amortise prefill across requests; DG's bidirectional
attention re-attends the full canvas on every denoising step and has no
bandwidth headroom to convert into throughput. The measured batch multipliers
are $\approx 80\times$ (AR-3B, vLLM; range 77--86$\times$) vs
$\dgBatchMult\times$ (DG, dInfer). At $\lambda=64$ ev/s the effective per-event
cost is $\approx 20$ ms (AR, $T_{\textsc{ar}}/80$) vs $\approx 321$ ms
(DG, $T_{\textsc{dg}}/\dgBatchMult$)---a $16\times$ AR advantage under load.
(This is the measured correlation; the causal mechanism is not
isolated---EDR-004.)

\textit{(e) Long input.}
AR pays for the prompt once; DG pays for it on every denoising step. DG's
input-length slope is $\dgInputSlope$ ms/prompt-tok vs AR's
$\arInputSlope$---$2.6\times$ steeper---so DG's prefill cost overwhelms its
canvas advantage past a crossover at $L_{\mathrm{in}}\approx 18{,}000$ tokens
(Table~\ref{tab:longinput}).

\begin{table}[ht]
\centering
\caption{Input-length sweep ($L_{\mathrm{out}}{=}256$): DG wins below the crossover, AR above.}
\label{tab:longinput}
\small
\begin{tabular}{rrrl}
\toprule
$L_{\mathrm{in}}$ (tokens) & AR-3B (ms) & DG (ms) & Winner \\
\midrule
1{,}000  & 1705 & 420  & DG ($4.1\times$) \\
4{,}000  & 1846 & 789  & DG ($2.3\times$) \\
7{,}000  & 1987 & 1158 & DG ($1.7\times$) \\
18{,}000 & 2504 & 2511 & crossover \\
25{,}000 & 2833 & 3372 & AR ($1.2\times$) \\
\bottomrule
\end{tabular}
\end{table}

\textit{(f) Deployment footprint at the quality-matched pair.}
The compute-matched pair is not the deployment pair: AR-1.5B reaches the same
(even higher) 9-field F1 at $17\times$ less memory, $\approx 17\times$ faster
cold start, and comparable per-event energy (Table~\ref{tab:deploy}). No
deployment axis favours DG at short output.

\begin{table}[ht]
\centering
\caption{Quality-matched deployment pair. At forced $L{\geq}512$, DG's latency
advantage widens dramatically because AR-1.5B rambles past its natural output
length ($\sim$123 tokens) when forced to generate longer sequences.}
\label{tab:deploy}
\small
\begin{tabular}{lcc}
\toprule
 & AR-1.5B & DG-26B \\
\midrule
9-field F1 ($L{=}256$)     & 0.8628 & \dgNineF \\
Weights (bf16) & 3.0 GB & 52.0 GB \\
Cold start      & $\sim$1 s  & $\sim$17 s \\
Energy/event (forced $L{=}256$) & 5.1 J  & 4.2 J \\
\midrule
\multicolumn{3}{l}{\textit{Per-event latency at forced output length:}} \\
$L{=}256$ (forced) & 1007 ms & 378 ms ($2.7\times$) \\
$L{=}512$ (forced) & 2005 ms & 501 ms ($4.0\times$) \\
$L{=}1024$ (forced) & 4026 ms & 739 ms ($5.4\times$) \\
\bottomrule
\end{tabular}
\end{table}

At natural output length ($L\sim 123$ tokens, no forcing), AR-1.5B completes
in 413 ms---only 35 ms slower than DG. But when the application requires
$L\geq 512$ tokens of structured output (e.g.\ a multi-field incident
report), AR-1.5B must generate well past its natural length, and its latency
quadruples while DG's grows by only 32\%. The deployment footprint favours
AR; the per-event latency at long forced output favours DG.

\paragraph{Energy.}
At the compute-matched pair, DG uses $4.3\times$ less energy per event than
AR-3B (net of idle: 4.2\,J vs 22.3\,J). At the quality-matched pair, DG uses
$1.2\times$ less than AR-1.5B at forced $L{=}256$ (4.2\,J vs 5.1\,J), widening
to $1.9\times$ at $L{=}512$ and $2.6\times$ at $L{=}1024$. The energy
advantage tracks the latency advantage and does not independently favour DG.

\paragraph{Synthesis.}
The per-event advantage (a--c) is real but narrow: it holds only at batch-1,
short input, and $L\ge 176$, and it does not survive the three deployment
comparisons at short output. AR wins on batched throughput ($80\times$ vs
$\dgBatchMult\times$ multiplier), on long input ($2.6\times$ steeper DG slope,
crossover $\approx 18$k tokens), and on footprint (the quality-matched AR
model is $17\times$ smaller). However, when the application requires
\emph{forced} output length $L\ge 512$---as in multi-field structured report
generation---AR-1.5B's latency quadruples (to 2005\,ms) while DG's grows by
only 32\% (to 501\,ms), a $4\times$ DG advantage. The cost shapes produce a
genuine per-event crossover; the deployment envelope closes it at short
output, but DG's advantage widens at long forced output.

\paragraph{When DM is the right choice.}
The decision procedure identifies a narrow but real niche. DM is the correct
architectural choice when \emph{all five} conditions hold:
\begin{enumerate}[noitemsep]
\item Output $\ge 512$ tokens (forced), where DG is $4\times$ faster than
the quality-matched AR-1.5B;
\item Batch-1 by requirement (privacy isolation, interactivity, or cold-path
analysis), not by low load;
\item Complete output required (schema validation), because DM degrades
\emph{quality} under deadline pressure while AR degrades \emph{completeness};
\item Quality floor requires $>$1.5B parameters (AR-1.5B not good enough);
\item 52\,GB footprint acceptable (data center, not edge).
\end{enumerate}
Additionally, DM offers a per-request quality knob ($K$): at $K{=}4$ it meets
a 400\,ms p99 SLA (measured: 394.3\,ms) with F1$=$0.79, while AR has no
equivalent runtime knob---only the discrete choice of deploying a smaller
model.
